\documentclass[11pt,a4paper]{article}

\usepackage[utf8]{inputenc}
\usepackage[T1]{fontenc}
\usepackage[brazilian]{babel}
\usepackage{lmodern}
\usepackage[margin=2.5cm]{geometry}
\usepackage{amsmath,amssymb}
\usepackage{xcolor}
\usepackage[most]{tcolorbox}
\usepackage{enumitem}
\usepackage{hyperref}
\usepackage{fancyhdr}
\usepackage{titlesec}
\usepackage{microtype}
\usepackage{array}
\usepackage{booktabs}
\usepackage{listings}

\definecolor{deitelblue}{RGB}{31,73,125}
\definecolor{deitelorange}{RGB}{198,89,17}
\definecolor{boapratcolor}{RGB}{0,112,60}
\definecolor{errocolor}{RGB}{178,34,34}
\definecolor{engcolor}{RGB}{31,73,125}
\definecolor{desempcolor}{RGB}{198,89,17}
\definecolor{portabcolor}{RGB}{102,46,145}
\definecolor{codebg}{RGB}{248,248,248}
\definecolor{codekeyword}{RGB}{0,0,180}
\definecolor{codestring}{RGB}{163,21,21}
\definecolor{codecomment}{RGB}{0,128,0}

\lstdefinestyle{cppcode}{
  language=C++,
  basicstyle=\ttfamily\small,
  keywordstyle=\color{codekeyword}\bfseries,
  stringstyle=\color{codestring},
  commentstyle=\color{codecomment}\itshape,
  numbers=left, numberstyle=\tiny\color{gray}, numbersep=8pt,
  backgroundcolor=\color{codebg}, frame=single, framesep=4pt,
  rulecolor=\color{deitelblue!50}, breaklines=true,
  showstringspaces=false, tabsize=4,
  morekeywords={string, size_t, cin, cout, endl, inline, template, typename,
                bool, true, false, nullptr, override, final, public, private,
                protected, explicit, friend, virtual, this},
  literate={ã}{{\~a}}1 {á}{{\'a}}1 {é}{{\'e}}1 {ê}{{\^e}}1 {í}{{\'i}}1
           {ó}{{\'o}}1 {ô}{{\^o}}1 {õ}{{\~o}}1 {ú}{{\'u}}1 {ç}{{\c{c}}}1
}

\hypersetup{colorlinks=true, linkcolor=deitelblue, urlcolor=deitelblue,
  pdftitle={C: Como Programar - Cap. 17 - Autorrevisao}}

\pagestyle{fancy}
\fancyhf{}
\fancyhead[L]{\small\textit{C: Como Programar} --- Cap.~17}
\fancyhead[R]{\small Exerc\'icios de Autorrevis\~ao}
\fancyfoot[C]{\small\thepage}
\renewcommand{\headrulewidth}{0.4pt}

\titleformat{\section}{\Large\bfseries\color{deitelblue}}{\thesection}{1em}{}
\titleformat{\subsection}{\large\bfseries\color{deitelblue}}{\thesubsection}{1em}{}

\newtcolorbox{boaprat}{enhanced, breakable, colback=boapratcolor!8, colframe=boapratcolor,
  fonttitle=\bfseries, title={\textbf{Boa Pr\'atica de Programa\c{c}\~ao}},
  left=8pt, right=8pt, top=4pt, bottom=4pt}
\newtcolorbox{errocomum}{enhanced, breakable, colback=errocolor!8, colframe=errocolor,
  fonttitle=\bfseries, title={\textbf{Erro Comum de Programa\c{c}\~ao}},
  left=8pt, right=8pt, top=4pt, bottom=4pt}
\newtcolorbox{obseng}{enhanced, breakable, colback=engcolor!8, colframe=engcolor,
  fonttitle=\bfseries, title={\textbf{Observa\c{c}\~ao de Engenharia de Software}},
  left=8pt, right=8pt, top=4pt, bottom=4pt}
\newtcolorbox{dicadesemp}{enhanced, breakable, colback=desempcolor!8, colframe=desempcolor,
  fonttitle=\bfseries, title={\textbf{Dica de Desempenho}},
  left=8pt, right=8pt, top=4pt, bottom=4pt}
\newtcolorbox{resposta}{enhanced, breakable, colback=deitelblue!5, colframe=deitelblue!70,
  boxrule=0.6pt, left=8pt, right=8pt, top=4pt, bottom=4pt}

\title{
  \vspace{-1cm}
  {\Huge\bfseries\color{deitelblue} Cap\'itulo 17}\\[0.3em]
  {\Large\bfseries Classes: uma vis\~ao mais detalhada, parte 1}\\[0.8em]
  {\large\itshape Exerc\'icios de Autorrevis\~ao --- Resolu\c{c}\~ao Did\'atica}
}
\author{}
\date{}

\begin{document}
\maketitle
\thispagestyle{empty}
\vspace{-2em}

\begin{center}
\rule{0.85\textwidth}{0.5pt}\\[0.4em]
\textit{``No Cap\'itulo 16 dimos o passo do procedural para o orientado a objetos. Agora aprofundamos: separa\c{c}\~ao expl\'icita entre interface (\texttt{.h}) e implementa\c{c}\~ao (\texttt{.cpp}), construtores e destrutores, atribui\c{c}\~ao default de objetos, e o operador bin\'ario \texttt{::} em todo o seu vigor. Os exerc\'icios de autorrevisão deste cap\'itulo s\~ao curtos --- apenas dois ---, mas concentrados em pontos sutis que costumam confundir o iniciante.''}\\[0.4em]
\rule{0.85\textwidth}{0.5pt}
\end{center}

\vspace{1em}

% ============================================================
\section{Exerc\'icio 17.1 --- Operadores de acesso e visibilidade}
% ============================================================

\textit{Preencha os espa\c{c}os em cada uma das senten\c{c}as (a)--(d).}

\subsection*{(a) Operadores que acessam membros de classe}

\begin{resposta}\textbf{Resposta:} operador \textbf{ponto} (\texttt{.}) em conjunto com o nome de um objeto ou refer\^encia; operador \textbf{seta} (\texttt{->}) em conjunto com um ponteiro para um objeto.\end{resposta}

H\'a duas situa\c{c}\~oes distintas em que precisamos acessar um membro de uma classe a partir de um objeto:

\begin{enumerate}[leftmargin=*,itemsep=0pt]
  \item \textbf{Quando temos o objeto diretamente (ou uma refer\^encia a ele):} usa-se o operador ponto.
    \begin{lstlisting}[style=cppcode,numbers=none]
Conta c(100);          // c e um objeto
c.credit(50);          // acesso via "."
Conta &r = c;          // r e uma referencia
r.credit(20);          // tambem "." (referencia e como apelido)
    \end{lstlisting}
  \item \textbf{Quando temos um ponteiro para o objeto:} usa-se o operador seta.
    \begin{lstlisting}[style=cppcode,numbers=none]
Conta *p = &c;         // p e ponteiro
p->credit(50);         // acesso via "->"
// equivalente a (*p).credit(50)
    \end{lstlisting}
\end{enumerate}

\begin{obseng}
O operador \texttt{->} \'e uma abrevia\c{c}\~ao did\'atica e sint\'aticamente conveniente para \texttt{(*p).}~. Os par\^enteses em \texttt{(*p).credit(50)} s\~ao indispens\'aveis em raz\~ao das precedencias: sem eles, \texttt{*p.credit(50)} seria interpretado como \texttt{*(p.credit(50))}, que n\~ao tem sentido (\texttt{p} \'e ponteiro, n\~ao tem membros diretamente). Por isso o operador \texttt{->} foi criado: torna o c\'odigo leg\'ivel e elimina o risco de erro de prec\^edencia.
\end{obseng}

\subsection*{(b) Membros vis\'iveis somente dentro da classe}

\begin{resposta}\textbf{Resposta:} \texttt{private}.\end{resposta}

Os membros \texttt{private} formam a \emph{implementa\c{c}\~ao} oculta da classe. Apenas as fun\c{c}\~oes-membro da pr\'opria classe e suas \emph{amigas} (\emph{friends}) podem acess\'a-los. Tipicamente, os dados-membro v\~ao para a se\c{c}\~ao \texttt{private}, e qualquer fun\c{c}\~ao auxiliar que n\~ao deve ser exposta aos clientes tamb\'em vai para l\'a.

\subsection*{(c) Membros acess\'iveis onde quer que o objeto esteja no escopo}

\begin{resposta}\textbf{Resposta:} \texttt{public}.\end{resposta}

Os membros \texttt{public} formam a \emph{interface} da classe --- aquilo que os clientes podem usar. Tipicamente, contem os construtores, o destrutor, e as fun\c{c}\~oes-membro que oferecem servi\c{c}os aos clientes.

\begin{boaprat}
A regra de ouro do encapsulamento: \emph{dados-membro em \texttt{private}, fun\c{c}\~oes-membro que oferecem servi\c{c}os em \texttt{public}}. Exce\c{c}\~oes existem (constantes p\'ublicas, classes de pol\'iticas), mas s\~ao excecionais. Quando um dado-membro precisa ser exposto, fa\c{c}a-o atrav\'es de \emph{getters}/\emph{setters} ---  como vimos no Cap\'itulo 16.
\end{boaprat}

\subsection*{(d) Forma de atribuir um objeto a outro da mesma classe}

\begin{resposta}\textbf{Resposta:} a \textbf{atribui\c{c}\~ao de membros default} (\emph{default memberwise assignment}), realizada pelo operador de atribui\c{c}\~ao \texttt{=}.\end{resposta}

C++ gera, automaticamente, um operador de atribui\c{c}\~ao para cada classe que voc\^e define. Esse operador padr\~ao copia, um a um, todos os dados-membro do objeto da direita para o objeto da esquerda:

\begin{lstlisting}[style=cppcode,numbers=none]
Conta c1(100);
Conta c2;
c2 = c1;       // atribui c1 a c2: copia cada dado-membro
\end{lstlisting}

\begin{errocomum}
A atribui\c{c}\~ao default funciona maravilhosamente para classes simples (cujos dados-membro s\~ao todos primitivos ou s\~ao tipos com semantica adequada de c\'opia). Mas, para classes que gerenciam recursos --- ponteiros para mem\'oria din\^amica, descritores de arquivo, conex\~oes de rede ---, a atribui\c{c}\~ao default \'e \emph{tipicamente perigosa}: dois objetos passam a apontar para o mesmo recurso, e o primeiro a ser destru\'ido o libera, deixando o outro com um ponteiro pendente. Para essas classes \'e preciso definir \emph{operador de atribui\c{c}\~ao personalizado}; tema do Cap\'itulo~19 (\emph{Big Three} ou, em C++ moderno, \emph{Rule of Five}).
\end{errocomum}

% ============================================================
\section{Exerc\'icio 17.2 --- Ca\c{c}a aos erros}
% ============================================================

\textit{Encontre o(s) erro(s) em cada item e explique como corrigir.}

\subsection*{(a) Destrutor com tipo de retorno e par\^ametro}

\textbf{Segmento original:}
\begin{lstlisting}[style=cppcode,numbers=none]
void ~Time( int );
\end{lstlisting}

\begin{resposta}
\textbf{Resposta oficial:} destrutores n\~ao podem retornar valores nem ter par\^ametros. \textbf{Corre\c{c}\~ao:} remover o tipo de retorno \texttt{void} e o par\^ametro \texttt{int}.
\end{resposta}

\textbf{Vers\~ao correta:}
\begin{lstlisting}[style=cppcode,numbers=none]
~Time();
\end{lstlisting}

\subsection*{Coment\'ario did\'atico}

A regra \'e mais r\'igida e abrange tamb\'em construtores. Em C++:

\begin{itemize}[leftmargin=*,itemsep=0pt]
  \item \textbf{Construtor:} pode ter par\^ametros (e sobrecargas), \emph{n\~ao} pode ter tipo de retorno (nem \texttt{void}); seu nome \'e o nome da classe.
  \item \textbf{Destrutor:} \emph{n\~ao} pode ter par\^ametros, \emph{n\~ao} pode ter tipo de retorno; seu nome \'e til (\texttt{\~{}}) seguido do nome da classe; toda classe pode ter no m\'aximo \emph{um} destrutor (n\~ao h\'a sobrecarga de destrutor).
\end{itemize}

\begin{obseng}
A justificativa para essas restri\c{c}\~oes \'e profunda: construtores e destrutores n\~ao s\~ao ``chamados'' como fun\c{c}\~oes comuns; eles s\~ao \emph{invocados implicitamente} pelo compilador em momentos espec\'ificos (cria\c{c}\~ao e destrui\c{c}\~ao do objeto). O compilador n\~ao tem nenhum lugar onde colocar o valor de retorno do destrutor --- a quem ele entregaria? Nem por que escolheria um par\^ametro espec\'ifico no momento de destruir um objeto. Da\'i a proibi\c{c}\~ao.
\end{obseng}

\subsection*{(b) Inicializa\c{c}\~ao expl\'icita na defini\c{c}\~ao da classe}

\textbf{Segmento original:}
\begin{lstlisting}[style=cppcode,numbers=none]
class Time
{
public:
    // prototipos de funcao
private:
    int hour = 0;
    int minute = 0;
    int second = 0;
};
\end{lstlisting}

\begin{resposta}
\textbf{Resposta oficial:} os membros n\~ao podem ser inicializados explicitamente na defini\c{c}\~ao da classe. \textbf{Corre\c{c}\~ao:} remover a inicializa\c{c}\~ao expl\'icita; inicializar os dados-membro em um construtor.
\end{resposta}

\textbf{Vers\~ao correta:}
\begin{lstlisting}[style=cppcode,numbers=none]
class Time {
public:
    Time();              // construtor que zera tudo
private:
    int hour;
    int minute;
    int second;
};

Time::Time() : hour(0), minute(0), second(0) {}
\end{lstlisting}

\subsection*{Coment\'ario importante sobre vers\~oes do C++}

\begin{obseng}
\textbf{Nota hist\'orica importante:} a resposta oficial reflete o padr\~ao C++ at\'e a vers\~ao C++03 (e o livro foi escrito em 2011). A partir do \textbf{C++11}, a linguagem incorporou um recurso chamado \emph{default member initializer} (inicializador de membro padr\~ao) que permite, sim, escrever \texttt{int hour = 0;} dentro da defini\c{c}\~ao da classe! Em C++11 e posteriores, o segmento original \emph{\'e v\'alido}. Mas, no contexto pedag\'ogico do livro (que adota C++03 nas primeiras edi\c{c}\~oes), a resposta oficial estava correta. Adote o seguinte: \emph{na vers\~ao did\'atica do livro, inicialize em construtor}.
\end{obseng}

\subsection*{(c) Construtor com tipo de retorno}

\textbf{Segmento original:}
\begin{lstlisting}[style=cppcode,numbers=none]
int Employee( string, string );
\end{lstlisting}

\begin{resposta}
\textbf{Resposta oficial:} construtores n\~ao podem retornar valores. \textbf{Corre\c{c}\~ao:} remover o tipo de retorno \texttt{int}.
\end{resposta}

\textbf{Vers\~ao correta:}
\begin{lstlisting}[style=cppcode,numbers=none]
Employee( string, string );
\end{lstlisting}

\subsection*{Por que essa regra?}

Imagine que constru\c{c}\~oes pudessem retornar valor:

\begin{lstlisting}[style=cppcode,numbers=none]
Employee e("Maria", "Silva");   // o que faria com o "valor de retorno"?
\end{lstlisting}

\noindent A express\~ao \texttt{Employee e("Maria", "Silva")} \emph{\'e} a inicializa\c{c}\~ao de \texttt{e} --- ela n\~ao \'e uma express\~ao que produza um valor consumivel separadamente. O compilador n\~ao teria onde colocar um eventual valor retornado.

\begin{errocomum}
N\~ao confunda ``construtor retorna o objeto'' com ``construtor tem tipo de retorno''. Construtores efetivamente \emph{produzem} um objeto novo --- mas esse objeto \'e o que est\'a sendo declarado, n\~ao um valor de retorno separado. Em C++ moderno, fun\c{c}\~oes \emph{de f\'abrica} (\emph{factory functions}) podem retornar objetos por valor; mas elas n\~ao s\~ao construtores --- s\~ao fun\c{c}\~oes comuns que internamente chamam um construtor.
\end{errocomum}

% ============================================================
\section*{S\'intese conceitual}
% ============================================================

Os dois exerc\'icios de autorrevisão do Cap\'itulo 17 testam, em conjunto, quatro pontos essenciais:

\begin{itemize}[leftmargin=*,itemsep=0pt]
  \item \textbf{Acesso a membros:} \texttt{.} para objeto/refer\^encia, \texttt{->} para ponteiro.
  \item \textbf{Visibilidade:} \texttt{private} encapsula, \texttt{public} exp\~oe a interface.
  \item \textbf{Atribui\c{c}\~ao default:} copia membros um a um; suficiente para tipos simples, perigosa para gerenciamento de recursos.
  \item \textbf{Regras sint\'aticas de construtor e destrutor:} sem tipo de retorno; destrutor sem par\^ametros e \'unico.
\end{itemize}

Esses quatro pontos s\~ao a base sobre a qual o cap\'itulo constr\'oi as t\'ecnicas mais avan\c{c}adas: separa\c{c}\~ao interface/implementa\c{c}\~ao em arquivos distintos, construtores com listas de inicializa\c{c}\~ao, valida\c{c}\~ao centralizada e classes que mant\^em \emph{estado coerente} mesmo diante de mudan\c{c}as adversas. Aplicaremos tudo isso nos exerc\'icios de programa\c{c}\~ao 17.4 a 17.15.

\vspace{1em}
\begin{center}\rule{0.5\textwidth}{0.4pt}\end{center}

\end{document}
